\section{Related Work}

\subsection{Iterative Code Repair}

Self-Debug~\cite{chen2023selfdebug} introduced execution-feedback loops for LLM code generation, using test results and stack traces to guide iterative refinement. Self-Refine~\cite{madaan2023selfrefine} generalized this to broader refinement tasks using LLM self-feedback. These methods establish execution feedback as the baseline for repair loops.

LDB (LLM Debugger)~\cite{zhong2024ldb} extended this paradigm with block-by-block verification, improving localization of errors. However, all these approaches rely on execution traces---they know \emph{where} code fails but not \emph{why}.

\subsection{Static Analysis for LLMs}

Blyth et al.~\cite{blyth2025static} represent the closest prior work. They integrated static analysis feedback (pylint, bandit) into LLM pipelines and demonstrated improvements in code quality:
\begin{itemize}
    \item Security issues: 40\% $\rightarrow$ 13\%
    \item Readability problems: 80\% $\rightarrow$ 11\%
\end{itemize}

Critically, they did not measure functional correctness (pass@k). Their contribution establishes that LLMs can interpret static warnings---but whether this improves benchmark performance remains unknown.

\subsection{The Gap We Address}

\begin{table}[h]
\centering
\caption{Comparison of feedback approaches in prior work.}
\label{tab:prior-work}
\begin{tabular}{@{}llll@{}}
\toprule
Prior Work & Feedback Type & Metric & Pass@k? \\
\midrule
Self-Debug & Execution & Functional correctness & \checkmark \\
Self-Refine & LLM self-feedback & Task completion & \checkmark \\
Blyth et al. & Static analysis & Quality (security, readability) & $\times$ \\
\textbf{This work} & Static + Execution & Methodology baseline & Pending \\
\bottomrule
\end{tabular}
\end{table}

The research gap is clear: no study measures static analysis impact on pass@k. We provide the methodology and baseline to enable this evaluation.

\subsection{Benchmarks}

HumanEval~\cite{chen2021humaneval} (164 problems) and MBPP~\cite{austin2021mbpp} (974 problems) are standard benchmarks for code generation. Pass@k metrics measure functional correctness---the fraction of problems solved within k attempts. We focus on HumanEval as the more widely-used benchmark, with MBPP extension as future work.
